%% kapa: exam=1 round=27 year=2016 period=1 form=A subject=economics subject_ko=경제학원론 seq=28 number=68
%% source: Q-Net 감정평가사 기출문제 (cst00309) artlSeq=5169402 file=extracted/1008/5169402/제27회 감정평가사 1차 문제/제27회 감정평가사 1교시문제(A형).pdf
%% conversion: kordoc 4.24.51 -> tools/gen1.py (auto); review=proofread-sonnet; flags=none
\begin{question}{68}
폐쇄경제인 B국의 총생산함수가 $Y = AK^{\alpha}L^{1-\alpha}$이다. 전염병으로 인하여 $L$이 갑자기 감소함으로써 발생하는 변화에 관한 설명으로 옳은 것을 모두 고른 것은? (단, $A$와 $K$는 일정하며, $Y$는 총생산, $\alpha$는 0과 1 사이의 상수값, $A$는 총요소생산성, $K$는 자본량, $L$은 노동량으로 인구와 같다.)
\begin{framedtext}
ㄱ. 총생산은 증가할 것이다. \par
ㄴ. 1인당 생산은 증가할 것이다. \par
ㄷ. 자본의 한계생산은 증가할 것이다. \par
ㄹ. 노동의 한계생산은 증가할 것이다. \par
\end{framedtext}
\begin{choices}
  \item ㄱ, ㄴ
  \item ㄱ, ㄷ
  \item ㄴ, ㄷ
  \item ㄴ, ㄹ
  \item ㄷ, ㄹ
\end{choices}
\end{question}
\answer{4}
